Le diagnostic des tumeurs cérébrales repose sur la distinction entre gliome, méningiome
et tumeur hypophysaire à partir d'images IRM --- une tâche essentielle pour la planification
thérapeutique, mais chronophage lorsqu'elle est réalisée manuellement. Ce rapport présente
un pipeline d'apprentissage automatique en deux phases conçu pour classifier des IRM
de tumeurs cérébrales à l'aide de descripteurs artisanaux optimisés. En Phase~1, nous
effectuons une recherche de paramètres guidée par la séparabilité sur \NPhaseOneConfigs{}
configurations de quatre familles de descripteurs (statistiques, motifs binaires locaux
--- LBP, transformée en ondelettes discrète --- DWT, et histogrammes de gradients orientés
--- HOG), en sélectionnant la meilleure configuration par famille sans entraîner de
classifieur. En Phase~2, les jeux de caractéristiques optimisés et leurs fusions sont
comparés à \NModels{} classifieurs ajustés par validation croisée à 5 plis, soit
\NRuns{} combinaisons évaluées. Le meilleur pipeline --- Full(opt) + SVM à noyau
polynomial --- atteint une précision de \BestAcc{} et un F1-macro de \BestFOne{}
(IC bootstrap 95\,\% : [\BestFOneCILow, \BestFOneCIHigh]) sur un jeu de données
équilibré de \NTotal{} échantillons. La validation statistique (test de McNemar,
test de Friedman et analyse de stabilité par graine aléatoire) confirme la robustesse
des résultats. L'ensemble du pipeline s'exécute sur CPU, sans apprentissage profond,
et produit des descripteurs interprétables adaptés à l'aide à la décision clinique.
